\documentclass[10pt,a4paper]{amsart}
\usepackage[margin=19mm]{geometry}
\usepackage{amsmath,amssymb,mathtools}
\usepackage[scr=rsfs,cal=euler]{mathalfa}
\usepackage{xfrac}
\usepackage[unicode,hidelinks,bookmarks=false]{hyperref}
\newcommand{\ie}{i.e.\ }
\newcommand{\cf}{cf.\ }
\newcommand{\resp}{respectively\ }
\newcommand{\Subsection}{\subsection}
\providecommand{\nobreakdash}{\nobreak\hbox{-}}
\setlength{\emergencystretch}{2em}
\multlinegap0pt
\usepackage{bm}
\usepackage{bbm}
\usepackage{dsfont}
\usepackage{xspace}
\usepackage{cancel}
\usepackage{mathtools}
\usepackage{amsthm}
\makeatletter
\@ifundefined{theorem}{\newtheorem{theorem}{Theorem}}{}
\@ifundefined{lemma}{\newtheorem{lemma}[theorem]{Lemma}}{}
\makeatother
\def\Xcal{\mathcal{X}}
\newcommand{\bD}{D_L}
\title[On the local-global principle for twists of abelian varietie]{On the local-global principle for twists of abelian varieties and Galois representations}
\author{Nirvana Coppola, Lorenzo La Porta, Matteo Longo}
\date{Source: arXiv:2602.17299v2 (CC BY 4.0)}
\begin{document}
\maketitle

\noindent\emph{Source and licence.} On the local-global principle for twists of abelian varieties and Galois representations. Version arXiv:2602.17299v2 (CC BY 4.0), distributed under the Creative Commons Attribution 4.0 International licence, as recorded in the arXiv licence metadata for this exact version and shown on the abstract page https://arxiv.org/abs/2602.17299v2. Full rights notice: licenses/arxiv-2602.17299-CC-BY-4.0.txt.\par\smallskip

\noindent\textit{Editorial scope.} Counted unit: the Lemma whose source label is \texttt{lemma: skew reps and H1} in the article (source0.tex lines 847--849, source label \texttt{lemma: skew reps and H1}), delivered together with the article's own complete original proof of it (source0.tex lines 850--884, one \texttt{proof} environment of 35 lines). No mathematics is written, repaired or re-derived: statement and proof are the article's own text line for line. Required context retained uncounted: none. Every definition, display and label of those lines is retained unchanged. Omitted from this package and not redistributed: 1--846, 885--1057. Nothing inside a retained excerpt is omitted or altered: every retained body is delivered exactly as the article's lines are written, so no omission receipt of any kind accompanies this package. Citations: the retained excerpts carry no citation command at all, so no bibliography entry of the article is transcribed and no reference is redistributed. Reference adaptation: none was needed and none was made; every reference inside the delivered unit resolves inside it, so no unresolved reference or citation is produced. Printed slips kept literal: no printed word, symbol, hypothesis or number of the counted statement or of its proof was altered. Layout and class: the article's own class is \texttt{amsart} with packages including geometry, fontenc, inputenc, lmodern, babel, csquotes, amsthm, amssymb, amsmath, amsfonts, mathrsfs, amscd, amsbsy, dsfont; the wrapper is the lane's pinned \texttt{amsart} editorial wrapper at 10pt on A4, which restates the theorem-like environments the retained text needs and adds the article's own macro definitions that the retained text needs (\texttt{\textbackslash Xcal}, \texttt{\textbackslash bD}), with extra wrapper packages (bm, bbm, dsfont, xspace, cancel, mathtools). The delivered numbering is the unit's own. Assets: no retained span contains a figure environment, a graphics-inclusion command, a PDF-page inclusion, a table or a TikZ picture; no image file is copied into this package, the package declares no asset dependency and no image file is redistributed with it. Licence: version arXiv:2602.17299v2 of this article is distributed under the Creative Commons Attribution 4.0 International licence, as recorded in the arXiv licence metadata for this exact version and shown on the abstract page https://arxiv.org/abs/2602.17299v2. A full rights notice is delivered at licenses/arxiv-2602.17299-CC-BY-4.0.txt. Characters: every retained excerpt is pure ASCII, so the delivered engine, XeLaTeX with its default Latin Modern text font, reports no missing character.
\begin{lemma}\label{lemma: skew reps and H1}
There is a natural bijection \(H^1(G, \bD^{\times, \chi}) \cong {}_{\bD^{\times, \chi} \backslash} \Xcal\).
\end{lemma}

\begin{proof}
We consider the map
\begin{align*}
    Z^1(G, \bD^{\times, \chi}) & \longrightarrow \Xcal,\\
    u & \longmapsto f_u \coloneqq \sum_{\sigma \in G} a_\sigma \sigma \longmapsto \left( a \mapsto \sum_{\sigma \in G} a_\sigma \sigma(a) u(\sigma)^{-1}\right).
\end{align*}
One can check that \(f_u\) is a morphism of rings using the fact that \(u\) is a \(1\)-cocycle. 
% namely,
% \begin{align*}
%         f_u((a_\sigma \sigma)(a_\tau \tau))(a) &= f_u((a_\sigma \sigma(a_\tau))(\sigma \tau))(a) = a_\sigma \sigma( a_\tau) (\sigma \tau)(a) u(\sigma \tau)^{-1} =\\
%          &= a_\sigma \sigma( a_\tau) (\sigma \tau)(a) 
%         (\sigma u(\tau))^{-1}u(\sigma)^{-1}=a_\sigma \sigma( a_\tau) \sigma (\tau(a) )
%        \sigma(u(\tau)^{-1}) u(\sigma)^{-1}\\
%         & = a_\sigma \sigma( a_\tau \tau(a) u(\tau)^{-1} ) u(\sigma)^{-1} = f_u(a_\sigma \sigma)(f_u(a_\tau \tau)(a)),
% \end{align*}
% for all \(a \in \bD\).
Moreover, if \(b \in \bD^{\times, \chi}\) and we replace \(u\) by \(\sigma \mapsto bu(\sigma)\sigma(b^{-1})\), we see immediately that \(f_u\) is replaced by \(b \star f_u\). Thus, \(u \mapsto f_u\) induces a natural map \(H^1(G, \bD^{\times, \chi}) \to {}_{\bD^{\times, \chi} \backslash} \Xcal\). Conversely, consider \(f \in \Xcal\) and set \(u_f(\sigma) = (f(\sigma)(1_{\bD^\chi}))^{-1}\). 
% Then, we have
% \[\begin{split}
%     u_f(\sigma \tau) &= (f(\sigma \tau)(1_{\bD}))^{-1}\\&= (f(\sigma)(f(\tau)(1_{\bD})))^{-1} 
%     \\
%     &= (\sigma(f(\tau)(1_{\bD})) f(\sigma)(1_{\bD}))^{-1} 
%     \\& = u_f(\sigma) \sigma(u_f(\tau)),
%     \end{split}
% \]
% so
One can check that \(u_f\) is a \(1\)-cocycle. Moreover, if we replace \(f\) by \(b\star f\), \(b \in \bD^{\times, \chi}\), we get \(u_{b \star f}(\sigma) = bu_f(\sigma)\sigma(b^{-1})\).
% \[
%     u_{b \star f}(\sigma) =
%     (f(\sigma)(b)b^{-1})^{-1} = (\sigma(b)f(\sigma)(1_{\bD})b^{-1})^{-1} =  b (f(\sigma)(1_{\bD}))^{-1} \sigma(b^{-1}) = bu_f(\sigma)\sigma(b^{-1}).
% \]
Therefore, the cohomology class \([u_f]\) only depends on the equivalence class \([f]\).
Thus, \([u] \mapsto [f_u]\) and \([f] \mapsto [u_f]\) are well-defined and inverse to each other.  
%(just notice that \(u_{f_u}(\sigma)=(f_u(\sigma)(1_{\bD}))^{-1}=(u(\sigma)^{-1})^{-1}=u(\sigma)\)).
\end{proof}

\end{document}
